% Generated draft: requires mathematical and source review.
\documentclass[11pt]{article}
\usepackage[margin=25mm]{geometry}
\usepackage{fontspec}
\setmainfont{Latin Modern Roman}
\usepackage{amsmath,amssymb,mathtools,unicode-math}
\setmathfont{Latin Modern Math}
\usepackage{enumitem,needspace,xurl,hyperref}
\setlength{\parindent}{0pt}
\setlength{\parskip}{5pt}
\setlength{\emergencystretch}{3em}
\newcommand{\N}{\mathbb{N}}
\newcommand{\Z}{\mathbb{Z}}
\newcommand{\Q}{\mathbb{Q}}
\newcommand{\R}{\mathbb{R}}
\newcommand{\C}{\mathbb{C}}
\newcommand{\F}{\mathbb{F}}
\newcommand{\A}{\mathbb{A}}
\newcommand{\PP}{\mathbb{P}}
\newcommand{\E}{\mathbb{E}}
\newcommand{\eps}{\varepsilon}
\newcommand{\dd}{\,\mathrm{d}}
\begin{document}
\section*{UnsolvedMath ID: 30003833}
\textbf{Source problem code:} catalogue entry
\subsection{Definitions and mathematical statement}
\begin{enumerate}
\item \textbf{Prescribed tame inertia.} Let \(G\) be a finite group, and let \(C_1,\ldots,C_{m(G)+r}\) be conjugacy classes of cyclic subgroups of \(G\). For a cyclic subgroup \(H\leq G\), its conjugacy class is \([H]_G=\{gHg^{-1}:g\in G\}\). The source assumes that these classes generate \(G\), but does not specify whether this means that representatives \(H_i\in C_i\) can be chosen with \(\langle H_1,\ldots,H_{m(G)+r}\rangle=G\), or that the subgroups in their combined conjugacy classes generate \(G\).

The question asks for a finite Galois extension \(L/\Q\), together with an isomorphism \(\iota:\operatorname{Gal}(L/\Q)\xrightarrow{\sim}G\), whose inertia groups have the prescribed conjugacy classes \(C_1,\ldots,C_{m(G)+r}\), and which is tamely ramified. For a rational prime \(\ell\) and a prime \(\mathfrak P\) of \(L\) above \(\ell\), the inertia subgroup \(I_{\mathfrak P}\) consists of decomposition-group elements acting trivially on the residue field; its conjugacy class is independent of the choice of \(\mathfrak P\). Tameness at finite primes means \(\ell\nmid e(\mathfrak P/\ell)\) for every such pair. The source does not define \(m(G)\), quantify or constrain \(r\), specify whether the list enumerates all ramified finite primes with multiplicity, or prescribe behavior at the real place. Consequently, the exact quantified inertia-realization assertion remains unresolved.

\item \textbf{Regular realization of finite \(p\)-groups.} A finite \(p\)-group is a finite group of order \(p^n\), where \(p\) is prime and \(n\geq0\). A finite Galois extension \(K/\Q(t)\), with \(t\) transcendental over \(\Q\), is regular over \(\Q\) if every element of \(K\) algebraic over \(\Q\) belongs to \(\Q\). Equivalently, in a common algebraic closure of \(\Q(t)\), one has \(K\cap\overline{\Q}=\Q\). The regular-realization question is
\[
\text{for every prime }p\text{ and every finite }p\text{-group }P,
\quad \exists\,K/\Q(t)\text{ finite Galois and regular over }\Q
\quad\text{such that}\quad
\operatorname{Gal}(K/\Q(t))\cong P.
\]
The original statement instead writes \(K\cap\Q=\Q\). Since \(\Q\subseteq K\), that equality alone imposes no regularity condition. The formulation above uses the standard meaning of the explicitly supplied word ``regular''; the discrepancy requires review.

\item \textbf{Surface-group completions and Galois groups.} Under the usual convention that the surface is closed, connected, and orientable, its genus-\(g\) fundamental group is
\[
\Sigma_g=\left\langle a_1,b_1,\ldots,a_g,b_g\ \middle|\ \prod_{i=1}^{g}[a_i,b_i]=1\right\rangle,
\qquad [a,b]=aba^{-1}b^{-1},\qquad g\geq2.
\]
This surface convention is not expressly specified in the record. Write
\[
\widehat{\Sigma_g}=\varprojlim_{N\trianglelefteq\Sigma_g,\ [\Sigma_g:N]<\infty}\Sigma_g/N,
\qquad
\widehat{\Sigma_g}^{\,(p)}=\varprojlim_{N\trianglelefteq\Sigma_g,\ \Sigma_g/N\text{ a finite }p\text{-group}}\Sigma_g/N.
\]
For a field \(F\), let \(G_F=\operatorname{Gal}(F^{\mathrm{sep}}/F)\) be its absolute Galois group, and let \(G_F^{(p)}\) denote its maximal pro-\(p\) quotient, the inverse limit of its continuous finite \(p\)-group quotients. All isomorphisms here are topological group isomorphisms.

For each \(g\geq2\), the original statement asks whether there is a field \(F\) with
\[
G_F\cong\widehat{\Sigma_g}.
\]
It also asks, with a prime \(p\) understood as a parameter, whether there is a field \(F\) satisfying
\[
\operatorname{char}(F)\ne p,
\qquad G_F^{(p)}\cong\widehat{\Sigma_g}^{\,(p)}.
\]
The record gives no characteristic condition for the first question and does not expressly quantify \(p\). Its cleaned statement describes the pro-\(p\) completion as an absolute Galois group ``in the corresponding characteristic.'' That wording does not specify the characteristic and may ask for \(G_F\cong\widehat{\Sigma_g}^{\,(p)}\), which is different from the maximal-pro-\(p\)-quotient formulation above. Both readings must remain visible pending source review.

\item \textbf{Decidability of open subgroups.} A profinite group \(P\) is projective if, for every continuous epimorphism \(\alpha:A\twoheadrightarrow B\) of profinite groups and every continuous homomorphism \(f:P\to B\), there is a continuous homomorphism \(\widetilde f:P\to A\) with \(\alpha\circ\widetilde f=f\). An open subgroup \(U\leq P\) is given the induced profinite topology.

The record does not specify the language, theory, presentation, or effectiveness convention intended by ``decidable.'' Conditional on a fixed intended meaning, write \(D(P)\) for that property. The cleaned statement asks whether
\[
\forall\,P\text{ projective profinite},\qquad
D(P)\ \Longrightarrow\
\bigl(\forall\,U\leq P\text{ open},\ D(U)\bigr).
\]
The original statement asks whether there exist a projective profinite group \(P\) and an open subgroup \(U\leq P\) such that \(D(P)\) holds and \(D(U)\) fails. With a fixed meaning of decidability, these are complementary yes-or-no questions about the same closure property. No definition of \(D\) is selected here.
\end{enumerate}
\subsection{Short English statement}
This record bundles four questions: realizing finite groups over \(\Q\) with prescribed tame inertia, realizing finite \(p\)-groups regularly over \(\Q(t)\), realizing surface-group completions as absolute Galois groups or maximal pro-\(p\) Galois quotients, and preservation of decidability under passage to open subgroups of projective profinite groups. Several parameters and conventions are missing, and the supplied versions differ on the pro-\(p\) target and on the polarity of the decidability question. These qualifications require human mathematical review.
\subsection{Sources}
\begin{itemize}
\item Supplied local UnsolvedMath problem registry, permanent ID \(30003833\), exact source problem code \texttt{catalogue entry}, title ``Inverse-Galois and Decidability Problems for Profinite Groups,'' in the collection ``Oberwolfach Reports Open Problems.'' No individual proposer is supplied; the proposed year is recorded as 2018.
\item Supplied report citation: \emph{Field Arithmetic (2019)}, Oberwolfach Reports, DOI 10.4171/owr/2018/25, pp.~1517--1577. Supplied source link: \url{https://doi.org/10.4171/owr/2018/25}. The background attributes the problems to workshop open-problem discussions documented in Oberwolfach Reports. The citation's 2019 label and the 2018 DOI component are retained as supplied; their relationship has not been checked.
\item Supplied literature source, labeled ``Oberwolfach Report 2019'': \url{https://ems.press/content/serial-article-files/46747}.
\item The supplied literature assessment was checked at the exact recorded date \texttt{2026-08-21}. This is supplied dated curation, not independently verified. It labels the bundle open and describes unresolved tame inverse-Galois, regular \(p\)-group realization, profinite, and decidability questions; it further states that known solvable and simple-group cases do not settle the general claims. The supplied background explicitly calls this dated literature triage derived from OWR curation and requests independent verification of time-sensitive status and cited claims.
\item The registry supplies status \texttt{open} and statement status \texttt{corrected\_verified}, with an explanation that notation, typography, and sentence structure were repaired while mathematical content was preserved. These are source-record labels, not independent verification of the formulations or their mathematical status. The authoritative report citation is supplied, but its contents have not been independently inspected for this draft.
\end{itemize}
\end{document}
